\section*{Propósito y alcance}
Las siguientes propuestas buscan incorporar tecnología, conocimiento y capital a la actividad productiva. Los referentes internacionales muestran instrumentos que pueden adaptarse a Colombia. Su existencia no demuestra que produzcan el mismo resultado en otro contexto. La selección de proyectos requiere evaluar costos, capacidades institucionales y beneficios para sus usuarios. Los programas colombianos constituyen antecedentes para decidir qué ampliar o ajustar.

\section{Concentrar la inversión pública en proyectos sectoriales con resultados verificables}
El PND debería priorizar una cartera de proyectos de inteligencia artificial (IA) con problemas definidos, financiación identificada y responsables de ejecución. Para cada proyecto, la entidad sectorial establecería una línea de base, el costo completo y los resultados esperados. El DNP coordinaría la priorización y la evaluación, mientras MinTIC apoyaría las condiciones técnicas comunes. El seguimiento mediría mejoras en tiempos, errores o costos, además de registrar las herramientas entregadas.

La respuesta británica al \emph{AI Opportunities Action Plan} ofrece un referente de responsabilidad explícita por la ejecución. El CONPES 4144 permite identificar compromisos y capacidades existentes en Colombia. La cartera debe seleccionar las intervenciones útiles para las prioridades del Gobierno y financiar también su operación y mantenimiento \cite{ukplan,conpes}.

\section{Vincular a desarrolladores de IA con empresas que necesitan soluciones}
Se propone una convocatoria de cofinanciación por etapas para acuerdos entre desarrolladores y empresas usuarias. El apoyo cubriría la formulación del problema, la preparación de datos, la adaptación de la solución y su validación en operación. La empresa usuaria aportaría recursos, conocimiento del proceso y personal para evaluar los resultados. Los desembolsos posteriores dependerían del cumplimiento de hitos.

BridgeAI, en el Reino Unido, conecta oferta tecnológica y demanda empresarial. AI Assist, en Canadá, combina asesoría y apoyo a empresas innovadoras. Para Colombia, la selección debería exigir un usuario comprometido, reglas sobre datos y propiedad de los desarrollos y una justificación del beneficio público. El indicador principal sería la mejora verificada para el usuario, junto con la continuidad comercial de la solución \cite{bridge,canada}.

\section{Extender la asistencia empresarial para que las mipymes adopten IA}
Fábricas de Productividad y Sostenibilidad puede incorporar una línea de diagnóstico, preparación e implementación de IA. La asistencia ayudaría a escoger herramientas, ordenar datos, adaptar procesos y formar a sus usuarios. El diagnóstico debe poder recomendar digitalización o automatización convencional cuando resuelva el problema a menor costo. El apoyo incluiría integración y una parte temporal del acceso a herramientas, con conocimiento previo de los costos posteriores al subsidio.

EDGE, de Singapur, y AI Assist, de Canadá, ofrecen referentes de financiación y acompañamiento técnico. Su adaptación requiere mecanismos de desembolso accesibles a empresas que no pueden anticipar todo el gasto. El seguimiento comprobaría uso después del apoyo y mejoras en productividad, comparando con empresas o procesos equivalentes \cite{fabricas,edge,canada}.

\section{Facilitar pruebas supervisadas en finanzas y seguros}
La Superintendencia Financiera debería organizar pruebas de aplicaciones de IA con productos, usuarios, plazos y criterios de salida definidos. Las capacidades existentes de innovación y supervisión servirían de base. Cada prueba debe permitir detener la aplicación, revisar decisiones y volver al procedimiento anterior. En detección de fraude, por ejemplo, se medirían tanto operaciones ilícitas identificadas como pagos legítimos bloqueados.

AI Live Testing, de la autoridad financiera británica, ofrece un referente de pruebas acompañadas. AI Pathfinder, de Singapur, aporta un ejemplo de proyectos delimitados con apoyo técnico y financiero. En Colombia, la participación pública debe concentrarse en capacidad de supervisión y aprendizaje compartido. Los subsidios a entidades financieras requieren justificar beneficios adicionales a los que recibe la entidad participante \cite{fca,mas}.

\section{Desplegar herramientas compartidas en la administración pública}
Se propone un servicio común para contratar y utilizar herramientas de IA en tareas administrativas recurrentes. Las entidades podrían compartir acceso, seguridad, evaluación de proveedores y componentes reutilizables, manteniendo permisos diferenciados sobre sus datos. Las aplicaciones iniciales incluirían búsqueda documental, clasificación de solicitudes y borradores sujetos a revisión del funcionario responsable.

Pair, en Singapur, muestra una experiencia de herramientas compartidas para el servicio público. Las guías de la Agencia Digital de Japón ofrecen un referente de contratación y uso. La evaluación colombiana debe incorporar licencias, integración, seguridad, revisión humana y mantenimiento. El tiempo liberado puede mejorar el servicio o reducir atrasos. Solo debe registrarse como ahorro fiscal cuando reduzca gasto o evite un costo que habría ocurrido \cite{pair,japon}.

\section{Formar para el estudio y el trabajo y apoyar la reconversión}
Se propone combinar formación docente y criterios de uso educativo con rutas laborales por ocupación. En el trabajo, los bonos de formación financiarían programas evaluados sobre tareas concretas, con práctica y comprobación de competencias. La reconversión incluiría orientación y conexión con empleadores. En educación, las pruebas deben medir aprendizaje y capacidad de verificar información, además del uso de herramientas.

Los programas educativos de Singapur y las guías docentes de Inglaterra ofrecen referentes de acompañamiento a públicos diferentes. El diseño colombiano debe distinguir alfabetización general, aplicaciones por oficio y formación técnica especializada. La renovación de proveedores dependería de resultados en competencias y aplicación efectiva. Los cambios en empleo o ingresos requieren una comparación que permita separar el efecto de la formación de otros factores \cite{educacion}.

\section{Especializar la adopción en servicios profesionales y administrativos}
Dentro de la asistencia a mipymes, se propone una línea para firmas contables, jurídicas, de consultoría y de servicios administrativos. Universidades y organizaciones profesionales ayudarían a identificar tareas, preparar casos de prueba y evaluar herramientas. La asistencia incorporaría trazabilidad de documentos, verificación de fuentes y protección de información, según las obligaciones de cada profesión.

Singapur ha planteado formación que combina conocimiento profesional y aplicaciones de IA. El plan británico de adopción en servicios profesionales aporta otro antecedente de diseño. Para Colombia, los apoyos compartidos permitirían distribuir costos de diagnóstico y validación entre varias firmas. La compra y operación deben contar con aportes de las beneficiarias y permitir elegir proveedores. Esta línea puede ejecutarse con el instrumento de la tercera propuesta \cite{profesiones}.

\section{Facilitar acceso a cómputo para proyectos que acrediten su necesidad}
Se proponen apoyos temporales y competitivos para contratar procesamiento en la nube o acceder a infraestructura compartida. Los solicitantes deberían acreditar datos utilizables, capacidad técnica, una aplicación viable y una restricción de cómputo que impida desarrollarla. El apoyo debe permitir comparar soluciones y comprobar resultados antes de asumir costos permanentes.

El AI Research Resource británico y el AI Compute Access Fund canadiense ofrecen referentes de acceso para investigación y empresas. En Colombia, la decisión entre comprar equipos y contratar servicios requiere comparar demanda, utilización y costo total. La inversión propia incluiría energía, enfriamiento, personal y reposición. Los servicios contratados exigirían evaluar costos recurrentes, portabilidad y seguridad. La continuidad dependería del beneficio y del costo de operar la aplicación \cite{computo}.

\section{Desarrollar alianzas estables de investigación con financiación empresarial}
Se proponen acuerdos plurianuales entre empresas, universidades y centros de investigación para resolver preguntas productivas. Un componente financiaría capacidades compartidas, como personal científico, laboratorios y mantenimiento. Otro financiaría proyectos específicos con presupuestos, hitos y derechos sobre resultados. La participación empresarial se comprobaría mediante aportes monetarios pagados y recursos en especie con valoración verificable.

Manufacturing USA e IACMI ofrecen referentes de capacidades tecnológicas compartidas y cofinanciación. Los acuerdos CRADA de los laboratorios estadounidenses muestran mecanismos para investigar con participantes externos. En Colombia, los convenios de cooperación científica y tecnológica pueden servir de base, según el objeto y el régimen aplicable. Los acuerdos deben definir uso de resultados, acceso al conocimiento, licenciamiento y condiciones de renovación. La financiación pública correspondería a beneficios colectivos identificables \cite{investigacion}.

\section{Mejorar la asignación del riesgo para ampliar la financiación mediante capital}
El eje de la reforma sería eliminar la garantía de rentabilidad mínima del ahorro pensional y la reserva de estabilización asociada, como recomendó la Misión del Mercado de Capitales. La comparación con pares puede incentivar carteras similares y desalentar estrategias diferenciadas. Se propone evaluar las inversiones por su contribución al riesgo y al resultado neto de una cartera diversificada, según el horizonte del afiliado. La pérdida de una inversión particular puede ser compatible con una gestión prudente. Eliminar la garantía modifica quién absorbe pérdidas y debe mejorar la relación entre riesgo y retorno para el ahorrador \cite{mision,iorp}.

Durante el primer año del PND, el Gobierno debería presentar una reforma del artículo 101 de la Ley 100 y del artículo 19, literal l, de la Ley 2381, junto con las normas concordantes. Debe armonizar la evaluación de los fondos generacionales y mantener la responsabilidad por negligencia, conflictos e incumplimientos. Como complemento, se propone reducir el piso de capital para las administradoras del componente complementario y graduar las exigencias por funciones, escala y riesgos. La revisión del artículo 2.43.1.2.1 del Decreto 2555 debe considerar también los mínimos de las licencias de origen. La transición requiere protección diferenciada para afiliados próximos al retiro \cite{normas}.

\section*{Ejecución y evaluación}
Cada intervención debe identificar responsables, población objetivo y costo completo antes de ampliar su cobertura. El seguimiento debe distinguir productos de resultados y registrar aportes públicos y privados efectivamente ejecutados. Las evaluaciones iniciales deben permitir decidir qué ampliar, ajustar o terminar. En financiación empresarial, es necesario separar aportes primarios de capital, compras de acciones existentes e inversiones de deuda dentro de fondos privados. El efecto sobre productividad y bienestar de los usuarios debe orientar la asignación de recursos.
